\clearpage
\VerdierPhysicalPageAnchor{0247}%
\begingroup\small\setlength{\parskip}{0pt}\setlength{\emergencystretch}{1.5em}
\setlength{\abovedisplayskip}{4pt plus 2pt minus 2pt}%
\setlength{\belowdisplayskip}{4pt plus 2pt minus 2pt}%

\VerdierSemanticAnchor{ch4d-4-6-5-resolution-type}
\noindent are resolutions of injective type (1.1.4).

\VerdierSemanticAnchor{ch4d-4-6-5-terminology}
Thus, what we call a \emph{Cartan--Eilenberg injective resolution of the
complex $Y$} is what these authors call simply an \emph{injective
resolution of the complex $Y$} ([1], Chapter~XVII), terminology that we
shall not use in this work. A Cartan--Eilenberg injective resolution of
$Y$ induces resolutions of injective type of the components of the
complex $Y$ and of the kernels of the differentials of $Y$
[\emph{loc.\ cit.}].

\VerdierSemanticAnchor{ch4d-4-6-5-finite-length}
We shall say that a Cartan--Eilenberg resolution
$\varepsilon\colon Y\longrightarrow I^{\bullet\bullet}$ has
\emph{finite length} if there is an integer $n_0$ such that
$I^{\bullet,n}=0$ for $n>n_0$. For a complex $Y$ to admit a
Cartan--Eilenberg injective resolution of finite length, it is necessary
and sufficient that the injective dimensions of the cohomology objects of
$Y$ and of the images of the differentials of $Y$ be bounded. Recall the
results established in \emph{loc.\ cit.}:

\VerdierSemanticAnchor{ch4d-4-6-6-a}
\medskip\noindent\textbf{Proposition 4.6.6.}\quad
\textit{a) When the category $\mathcal A$ has enough injective objects
(3.1.5), every simple complex $Y$ of $\mathcal A$ admits a
Cartan--Eilenberg injective resolution.}

\VerdierSemanticAnchor{ch4d-4-6-6-b}
\noindent\textit{b) Let
$f\colon X\longrightarrow Y$ be a morphism of simple complexes, and let
$\varepsilon\colon X\longrightarrow I^{\bullet\bullet}$ and
$\varepsilon'\colon Y\longrightarrow J^{\bullet\bullet}$ be two
Cartan--Eilenberg injective resolutions. There is a morphism of double
complexes}

\VerdierSemanticAnchor{ch4d-4-6-6-b-morphism}
\[
\textit{$g\colon I^{\bullet\bullet}\longrightarrow J^{\bullet\bullet}$}
\]
\textit{over $f$; i.e.\ such that the diagram of morphisms of double
complexes}

\VerdierSemanticAnchor{ch4d-4-6-6-b-diagram}
\[
\begin{tikzcd}[column sep=5.0em,row sep=3.0em]
X \arrow[r,"f"] \arrow[d,"\varepsilon"'] &
Y \arrow[d,"\varepsilon'"] \\
I^{\bullet\bullet} \arrow[r,"g"'] &
J^{\bullet\bullet}
\end{tikzcd}
\]
\textit{is commutative.}

\VerdierSemanticAnchor{ch4d-4-6-6-c}
\noindent\textit{c) Two morphisms
$g,g'\colon I^{\bullet\bullet}\longrightarrow J^{\bullet\bullet}$
over $f$ are homotopic (in the sense of homotopy of morphisms of double
complexes (Chapter~I, 2.5.1)).}

\VerdierSemanticAnchor{ch4d-4-6-6-proof}
For the proof, we refer to \emph{loc.\ cit.}

\VerdierSemanticAnchor{ch4d-4-6-7}
\medskip\noindent\textbf{4.6.7.}\quad
Let $Y$ be an object of $\operatorname{comp}(\mathcal A)$, and let
$Y\longrightarrow I^{\bullet\bullet}$ be a Cartan--Eilenberg injective
resolution of $Y$. The double complex $I^{\bullet\bullet}$ gives rise to
two spectral objects with values in $D(\mathcal A)$, according as it is
filtered by its first or its second degree (4.1.7). The spectral object
with values in $D(\mathcal A)$ defined by $I^{\bullet\bullet}$ filtered by
its second degree is called the \emph{Cartan--Eilenberg spectral object of
$Y$} (relative to the Cartan--Eilenberg resolution under consideration).
\endgroup
